\section{Threats to Validity}
\label{sec:threats}

\subsection{Scope of the conclusion}

We audit one published centrality-weighted optimizer and one adaptive scheme of
our own. The controls are defined for the general form \eqref{eq:master}, and
Proposition~\ref{prop:mean} holds for any score-proportional weighting, but the
empirical finding is about these two mechanisms. Whether other structural
statistics---closeness, eigenvector centrality, PageRank---or other graph
families behave the same way is untested here. We regard this as the principal
limitation: the instrument is general, the evidence is not yet.

Two considerations bear on how far the result plausibly extends. The argument
of Section~\ref{sec:discussion} turns on the graph being fixed before the
first evaluation and unrelated to the objective, which is true of every
Watts--Strogatz-based method we are aware of. And two independently designed
mechanisms failed their controls in the same way. Neither is evidence about
methods we did not run.

\subsection{Statistical power}

The primary contrast is run at the $51$ repetitions the CEC2017 protocol
prescribes; the remaining arms and the $D=10$ suite are run at $30$, which we
note as a limitation rather than a convention.

Per-function equivalence is nonetheless not established at the negligibility
margin, and additional runs would not fully establish it. At $51$ runs the
median $90\%$ half-width is $0.139$, below the $0.147$ margin, yet only $6$ of
$28$ functions are formally equivalent, because the criterion is
$|\hat{\delta}|+\text{half-width}<m$ rather than $\text{half-width}<m$. On the
median function, reaching the negligibility margin would require roughly $100$
runs; on $4$ of $28$ the point estimate itself exceeds $0.147$, so no sample
size would place them inside that band. Our equivalence claim is therefore made
at suite level, where the interval $[-0.038,+0.020]$ is contained in
$|\delta|<0.05$, and per function we claim only that effects lie within a small
margin ($26$ of $28$ at $|\delta|<0.33$) and that none survives correction as a
difference.

This is a real limit on the strength of the negative result, and we prefer to
state it than to rest on the difference tests alone. What the data exclude is a
systematic effect of the assignment across a suite; what they do not exclude is
a small effect on particular functions.

\subsection{Construct validity of the controls}

The permutation control assumes the score multiset is the quantity to hold
fixed. If a method's benefit came from the \emph{joint} distribution of scores
and degrees---high-degree nodes systematically receiving high scores---then
permutation would break that too, and attribute to assignment something that
belongs to a degree--score coupling. We regard this as a distinction without a
practical difference here, since such a coupling is exactly the kind of
structural information the claim asserts; but a reader who separates them
should read our result as covering both.

The policy control randomises selection while preserving the switching
schedule. It therefore tests whether selection is informed, not whether the
epoch length or the arm set are well chosen. A method whose benefit lay
entirely in its schedule would pass our control while making a weaker claim
than its authors intend.

\subsection{Implementation and benchmark choices}

All arms share one implementation, which removes a large source of variance but
means an implementation-level defect would affect every arm. We verified the
vectorized update against the reference path to $8.9\times10^{-16}$ per step
and to identical run distributions before using it. Population size is fixed at
$N=60$ throughout; the density sweep varies degree but not $N$, so interactions
between the two are unmeasured.

Two CEC2017 functions are absent, $F_{2}$ by the organisers' advice and
$F_{30}$ because the reference implementation we use omits it. The protocol's
fourth dimension, $D=100$, is not run; at several times the cost of $D=50$ it
was beyond the single machine available, and we note it as unfinished rather
than unnecessary. The primary contrast is run at all three dimensions we do
report; the density and adaptive arms are run at $D\in\{10,30\}$, since the
$7{,}350$-run sweep already covers density at $D\in\{10,30,50\}$.

\subsection{Reproducibility}

Every run is stored individually, keyed by arm, function, dimension and seed,
and every figure and table is produced by a script from those files. Seeds are
shared across arms so all comparisons are paired. Runs were executed on one
machine with a fixed software environment; we report the parallel throughput
measurements that determined the scheduling, since they affect wall-clock
figures but not results.
